\chapter{Literature Review}
\label{ch:literature}

This chapter surveys the intellectual lineage that informs the three pillars
of this thesis: optimal bet sizing under the Kelly Criterion, the
exploration-exploitation trade-off in Multi-Armed Bandits, and
risk-constrained portfolio management via CPPI.  We also situate our proposed
contributions---the Volatility-Augmented HMM and the Risk-Constrained Kelly
framework---within the existing literature.

% ---------------------------------------------------------------
\section{The Kelly Criterion: Origins and Controversy}
% ---------------------------------------------------------------

\subsection{Kelly (1956) and Asymptotic Growth Optimality}

The Kelly Criterion originates in a 1956 paper by \citet{kelly1956new}, a
researcher at Bell Labs, who drew an analogy between information-theoretic
channel capacity and optimal gambling strategies.  Kelly showed that an agent
maximizing the expected growth rate of wealth should maximize $\E[\log W_T]$,
and derived the closed-form fraction now bearing his name.  \citet{breiman1961}
subsequently supplied the rigorous measure-theoretic proof of asymptotic
dominance: any strategy that differs from Kelly on a set of positive measure
is eventually dominated almost surely.  This result is stated formally as
Theorem~\ref{thm:kelly_slln} in Chapter~\ref{ch:foundations}.

\subsection{The Samuelson--Kelly Debate}

Despite its theoretical appeal, Kelly betting has been criticized sharply,
most famously by the Nobel laureate \citet{samuelson1979}.  In a deliberately
monosyllabic 1979 paper, Samuelson argued that maximizing $\E[\log W]$ is
only justified if the investor's utility function is precisely logarithmic,
and that investors with power utility $U(W) = W^\gamma/\gamma$ ($\gamma \neq 0$)
should use \emph{fractional} Kelly strategies scaled by $\gamma$.

Samuelson's critique centers on two related points.  First, Kelly betting can
produce spectacular drawdowns over finite horizons even when the asymptotic
growth rate is maximal.  The median wealth under full Kelly can be far below
the expected wealth because the distribution of $W_T$ is log-normal with a
heavy right tail: expected wealth is driven by rare high-growth paths, while
the typical investor suffers significant downside.  Second, for an investor
who retires at a fixed horizon $T$, maximizing long-run growth is not the same
as maximizing the probability of reaching a target.  \citet{thorp2008} showed
that half-Kelly betting maximizes this probability in many practical settings,
even though it sacrifices asymptotic growth.

The tension between Samuelson's critique and Kelly's optimality result
motivates our fractional Kelly and CPPI experiments in
Chapter~\ref{ch:experiments}.  We empirically quantify the ``cost of
survival'': the foregone geometric growth rate when a drawdown floor is
imposed.

\subsection{Heavy Tails and the Kelly-Ruin Paradox}

Classical Kelly theory assumes finite-variance returns.  \citet{mandelbrot1963variation} and
\citet{fama1965behavior} documented that equity returns exhibit leptokurtosis---heavier
tails than the Gaussian.  For distributions in the domain of attraction of a
stable law with tail index $\alpha < 2$, the variance is infinite and the Kelly
fraction $f^* = \mu/\sigma^2$ computed from sample moments is systematically
overestimated.  In the limit, a finite-sample Gaussian Kelly fraction applied
to a Student-$t_3$ environment guarantees eventual ruin \citep{ziemba1986beat}.  This is the \emph{Kelly-Ruin Paradox}: the formula designed to prevent
ruin causes ruin when its Gaussian assumption is violated.  We reproduce and
quantify this paradox in Experiment~6.

% ---------------------------------------------------------------
\section{Multi-Armed Bandits: Exploration vs.\ Exploitation}
% ---------------------------------------------------------------

\subsection{The Robbins Problem and the Lai--Robbins Lower Bound}

The Multi-Armed Bandit problem was formalized by \citet{robbins1952} as a
sequential statistical decision problem.  \citet{lai1985asymptotically} established
the fundamental lower bound on regret: any consistent policy must suffer
expected regret of at least $\Omega(\log T)$.  This bound was achieved
constructively by UCB1 \citep{auer2002finite}, whose regret
guarantee (Theorem~\ref{thm:ucb_regret}) is logarithmic and tight.

\subsection{Thompson Sampling}

\citet{thompson1933likelihood} proposed probability matching---selecting actions with
probability proportional to their posterior probability of being optimal---over
eight decades before it was recognized as a competitive algorithm.  Rigorous
analysis was provided by \citet{agrawal2012analysis}, confirming asymptotic optimality (Theorem~\ref{thm:ts_regret}).

The combination of Thompson Sampling with Kelly bet sizing used in this thesis
is a novel contribution.  Standard bandit algorithms treat arm selection and
resource allocation as separate problems; our formulation couples them so that
Bayesian posterior uncertainty directly modulates the bet size, yielding an
implicit risk-aversion mechanism without any explicit risk constraint.

\subsection{Adversarial Bandits and EXP3}

\citet{auer2002nonstochastic} introduced EXP3 for the
adversarial bandit setting, where rewards may be chosen by an adversary with
full knowledge of the agent's strategy.  The $O(\sqrt{T K \ln K})$ minimax
regret bound of EXP3 (Theorem~\ref{thm:exp3_regret}) is qualitatively stronger
than UCB in nonstationary settings because it makes no stationarity assumption
whatsoever.  We use a fractional variant as a principled adversarial baseline,
noting the deviation from canonical EXP3 in Remark~\ref{rem:exp3_variant}.

% ---------------------------------------------------------------
\section{Regime-Switching Models and Change-Point Detection}
% ---------------------------------------------------------------

\subsection{Hidden Markov Models in Finance}

Hidden Markov Models were applied to financial time series by \citet{hamilton1989new},
who modeled GNP growth as a two-state Markov-switching process.  \citet{ang2002regime} documented that
equity returns are well described by regime-switching models with distinct bull
and bear states.  The parameters for our simulation environments
($\mu_{\text{bull}} = 0.08$, $\sigma_{\text{bull}} = 0.15$;
$\mu_{\text{bear}} = -0.10$, $\sigma_{\text{bear}} = 0.30$) are calibrated to
these empirical estimates.

\subsection{CUSUM Change-Point Detection}

\citet{page1954} introduced the Cumulative Sum (CUSUM) procedure for sequential
change-point detection.  \citet{moustakides1986optimal} established the minimax optimality of
CUSUM for minimizing expected detection delay subject to a false-alarm
constraint.  Our Vol-HMM augments the standard HMM filter with a CUSUM layer
precisely because CUSUM provides a rapid response to mean shifts that the
HMM's transition dynamics alone may miss during the first few post-switch
observations.  The combination---soft Bayesian filtering for steady-state
regime tracking plus hard CUSUM triggering for abrupt shifts---exploits the
complementary strengths of the two approaches.

% ---------------------------------------------------------------
\section{Risk-Constrained Portfolio Management}
% ---------------------------------------------------------------

\subsection{CPPI: Constant Proportion Portfolio Insurance}

\citet{black1992theory} formalized CPPI as a dynamic asset allocation strategy
that preserves a minimum wealth floor.  The key formula,
$\text{Exposure}_t = m \cdot (W_t - \text{Floor}_t)$, was shown to guarantee
(in continuous time) that wealth never falls below the floor.  In discrete time,
slippage through the floor is possible during large single-period drawdowns, a
limitation we document experimentally.

\citet{busseti2016risk} extended the CPPI idea into a convex optimization
framework, computing bet size via constrained log-growth maximization.  Our
implementation follows the spirit of their drawdown-constrained Kelly framework,
using the CPPI cushion formula to modulate the Kelly leverage dynamically.  The
formal guarantee is stated in Proposition~\ref{prop:cppi_bound} of
Chapter~\ref{ch:methodology}.

\subsection{Sharpe and Sortino Ratios}

\citet{sharpe1966mutual} proposed the reward-to-variability ratio as a risk-adjusted
performance measure; the annualized form is
$\mathcal{S} = (\bar{R}/\hat\sigma)\sqrt{252}$.  The Sortino ratio \citep{sortino1991downside} replaces total volatility with downside volatility,
making it more appropriate for strategies that explicitly bound drawdowns.
We compute both metrics in $O(T)$ via running-statistics recursions to avoid
storing full return histories in the simulation loop.

% ---------------------------------------------------------------
\section{Positioning of This Thesis}
% ---------------------------------------------------------------

This thesis sits at the intersection of three bodies of work: the Bayesian
bandit literature \citep{thompson1933likelihood, auer2002finite, auer2002nonstochastic}, the Kelly/growth-rate optimization
literature, and the regime-detection and risk-control literature \citep{hamilton1989new, page1954, black1992theory}.  The specific combination---a Bayesian filter over HMM states that
drives a Kelly-sized bet subject to a CPPI floor, evaluated under a unified
CRN simulation framework with empirical S\&P~500 validation---does not appear
to have been studied as a single integrated system in the prior literature.
The closest antecedents are \citet{thorp2008} on dynamic Kelly sizing
and \citet{ang2002regime} on regime-conditional portfolio allocation, but
neither combines all three components in the manner presented here.
